\documentclass[11pt]{article}
\usepackage{amsmath,amssymb}
\usepackage[margin=0.85in]{geometry}
\DeclareMathOperator{\tr}{tr}
\newcommand{\Mat}{\operatorname{M}_3}
\begin{document}
\begin{center}
  {\Large\bfseries A degree-eight central polynomial for \(3\times3\) matrices\par}
  \smallskip
  Disproof of Conjecture 00000009998
\end{center}

We exhibit a nonzero-valued homogeneous central polynomial of degree \(8\) on
\(\Mat(\mathbb Q)\). Since \(8<3^2\), its existence disproves the conjectured
minimum. Drensky and Kasparian previously established the existence of
degree-eight central polynomials for \(3\times 3\) matrices
\cite{drensky-kasparian}. The construction below includes a proof of centrality
and a concrete nonzero evaluation.

Let \(X,Y,C_1,C_2,C_3,C_4\) be matrix variables and put \(A=XY-YX\).
Introduce odd generators \(\theta_1,\ldots,\theta_4\) of an exterior algebra
over \(\mathbb Q\) and let \(U=\sum_{j=1}^4\theta_j C_j\). Define \(p\) as
the coefficient of \(\theta_1\theta_2\theta_3\theta_4\) in
\begin{equation}
\label{eq:expression}
 H=\sum_{a+b+c=4}U^a A U^b A U^c
       -2\bigl(UA^2U^3+UAU^2AU+U^3A^2U\bigr),
\end{equation}
where \(a,b,c\) are nonnegative integers. Equivalently, each degree-four
string of occurrences of \(U\) in \eqref{eq:expression} is replaced by the
alternating sum of the \(4!\) assignments of \(C_1,\ldots,C_4\) to those
positions. This makes \(p\) an ordinary noncommutative polynomial in its six
matrix variables. Each term after substitution uses four occurrences of
\(C_j\) and two commutators of degree two, for total degree eight.

\section*{Centrality}

For matrices with homogeneous exterior-algebra entries, the trace satisfies
\(\tr(RS)=(-1)^{\lvert R\rvert\lvert S\rvert}\tr(SR)\).
Put \(V=U^2\), \(T=UAU\), \(q=\tfrac12\tr(A^2)\),
\(b=\tr(AV)\), \(c=\tr(A^2V)\), and \(e=\tr(AV^2)\).
The matrices \(A,V,T\) have entries in the commutative even subalgebra. In
particular, Cayley--Hamilton applies there. The ordinary trace of a
commutator vanishes, so \(\tr(A)=0\). The signed trace identity gives
\[
 \begin{gathered}
   \tr(V)=\tr(V^2)=0,\qquad \tr(T)=-b,\qquad \tr(AT)=0,\\
   \tr(VT)=-e,\qquad \tr(AVT+ATV)=0.
 \end{gathered}
\]
For example, \(\tr(AT)=\tr((AU)(AU))\) changes sign when the two odd
factors are exchanged. For the final relation the two summands, up to an
even cyclic reordering, are the traces of the products of \(AU\) and
\(AU^3\) in opposite orders.

Write \(D=A^2V+AVA+VA^2\). Extracting the \(t\)-coefficient in the
Cayley--Hamilton identity for the traceless matrix \(A+tV\) gives
\[
    D=qV+bA+cI.
\]
To see the scalar term, use
\(\det(A+tV)=\frac13\tr((A+tV)^3)\) for a traceless size-three matrix.
Multiplying that Cayley--Hamilton identity by \(A+tV\) and extracting the
\(t^2\)-coefficient gives
\[
    Q=qV^2+b(AV+VA)+cV+eA,
\]
where \(Q\) is the sum of the six words with two copies each of \(A\) and
\(V\). Here the coefficient of \(t^2\) in \(\tfrac12\tr((A+tV)^2)\)
vanishes by \(\tr(V^2)=0\).

Let \(S\) be the sum of the six words that use \(A,V,T\) once each.
For \(M=xA+yV+zT\), extract the \(xyz\)-coefficient from
\[
 M^3=(\tr M)M^2-\tfrac12\bigl((\tr M)^2-\tr(M^2)\bigr)M+\det(M)I.
\]
The determinant contributes \(-\tr(T)\tr(AV)+\tr(AVT+ATV)=b^2\), by the
usual degree-three Newton identity. The other terms give
\[
    S=-b(AV+VA)+bT-eA+b^2I.
\]
In \eqref{eq:expression}, grouping the terms by which of \(a,b,c\) are
odd yields \(H=Q+S-UDU\). The terms with three even exponents give \(Q\).
The terms with two odd exponents adjacent around one occurrence of \(A\)
give \(S\); those with odd outer exponents have sign minus and give
\(-UDU\). Since \(UDU=qV^2+bT+cV\), it follows that \(H=b^2 I\).
Taking the exterior coefficient proves that all evaluations of \(p\) are
scalar matrices.

\section*{Nonzero evaluation}

With indices \(0,1,2\), write \(E_{ij}\) for the matrix unit at row
\(i\), column \(j\). Choose
\[
 X=E_{01},\quad Y=E_{10},\qquad
 (C_1,C_2,C_3,C_4)=(E_{01},E_{10},E_{02},E_{20}).
\]
Then \(A=\operatorname{diag}(1,-1,0)\). In \(\tr(AU^2)\), the coefficients
of \(\theta_1\theta_2\) and \(\theta_3\theta_4\) are \(2\) and \(1\),
respectively. All others vanish since their matrix commutators have zero
diagonal. Consequently the \(\theta_1\theta_2\theta_3\theta_4\)-coefficient
of \(b^2\) is \(4\), so this evaluation of \(p\) is \(4I\). Evaluating at
all zero matrices gives zero. Thus \(p\) is a genuine central polynomial
with nonconstant evaluation. Its degree is exactly eight, which suffices
to disprove the conjecture.

\section*{Exact coefficient verification}
The Lean project independently verifies centrality on diagonal trace-zero
matrices by calculating every matrix-unit coefficient. Here \(A\) is
treated as a separate traceless input before the substitution
\(A=XY-YX\). With \(A=\operatorname{diag}(a,b,-a-b)\), a nonzero ordered
product of four matrix units is a path through five vertices
\(i_0,\ldots,i_4\). Placing the two diagonal factors after \(r\) and
\(s\) units multiplies the path's output \(E_{i_0 i_4}\) by
\(A_{i_r i_r}A_{i_s i_s}\), where \(0\le r\le s\le4\).
The sign of the placement is minus when both \(r\) and \(s\) are odd
and plus otherwise, as in \eqref{eq:expression}.
The sign of the edge order is determined by its inversion parity in
row-major order. Repeated edges make no contribution by alternation.

The computation enumerates the three intermediate vertices for each
of the nine output entries and groups the contributions by ascending
edge tuple. It records the coefficients of \(a^2,ab,b^2\) as
integers. For every diagonal output it obtains the same table below.
The off-diagonal tables are empty. An absent tuple has zero
coefficient. The unit indices range from \(0\) to \(8\).
\[
\begin{array}{c|rrr}
\text{ascending edge tuple} & a^2 & ab & b^2\\ \hline
(1,2,3,6) & -4 & 2 & 2\\
(1,3,5,7) & 2 & 2 & -4\\
(2,5,6,7) & -4 & -10 & -4
\end{array}
\]
The table equality includes every four-element subset of the matrix
units. Alternation and multilinearity then extend it to all four matrix
arguments. The coefficient identities have integer coefficients and
therefore also hold over \(\overline{\mathbb Q}\). The verification
includes the commutator and the value \(4I\) at the stated witness.

For a general traceless input \(A\), let \(D_0=\operatorname{diag}(1,0,-1)\).
The discriminant of the characteristic polynomial of \(D_0+tA\) is a
polynomial in \(t\) and is nonzero at \(t=0\). For all but finitely many
rational \(t\), the matrix \(D_0+tA\) thus has distinct eigenvalues and
is conjugate over \(\overline{\mathbb Q}\) to a trace-zero diagonal
matrix. Simultaneous conjugation commutes with the polynomial evaluation,
so the off-diagonal entries and the differences of diagonal entries
vanish at all those \(t\). They are polynomials in \(t\). Their coefficients
of \(t^2\) are the corresponding entries and differences at \(A\), because
the expression is homogeneous quadratic in the traceless input.
These coefficients vanish, which extends the exact calculation to
arbitrary traceless inputs.

\begin{thebibliography}{1}
\bibitem{drensky-kasparian}
V. Drensky and A. Kasparian, \emph{A new central polynomial for
\(3\times3\) matrices}, Communications in Algebra \textbf{13} (1985),
745--752. DOI: \texttt{10.1080/00927878508823188}.
\end{thebibliography}
\end{document}
